% Options for packages loaded elsewhere
\PassOptionsToPackage{unicode}{hyperref}
\PassOptionsToPackage{hyphens}{url}
%
\documentclass[
  english,
  twoside]{article}
\usepackage{lmodern}
\usepackage{amssymb,amsmath}
\usepackage{ifxetex,ifluatex}
\ifnum 0\ifxetex 1\fi\ifluatex 1\fi=0 % if pdftex
  \usepackage[T1]{fontenc}
  \usepackage[utf8]{inputenc}
  \usepackage{textcomp} % provide euro and other symbols
\else % if luatex or xetex
  \usepackage{unicode-math}
  \defaultfontfeatures{Scale=MatchLowercase}
  \defaultfontfeatures[\rmfamily]{Ligatures=TeX,Scale=1}
\fi
% Use upquote if available, for straight quotes in verbatim environments
\IfFileExists{upquote.sty}{\usepackage{upquote}}{}
\IfFileExists{microtype.sty}{% use microtype if available
  \usepackage[]{microtype}
  \UseMicrotypeSet[protrusion]{basicmath} % disable protrusion for tt fonts
}{}
\makeatletter
\@ifundefined{KOMAClassName}{% if non-KOMA class
  \IfFileExists{parskip.sty}{%
    \usepackage{parskip}
  }{% else
    \setlength{\parindent}{0pt}
    \setlength{\parskip}{6pt plus 2pt minus 1pt}}
}{% if KOMA class
  \KOMAoptions{parskip=half}}
\makeatother
\usepackage{xcolor}
\IfFileExists{xurl.sty}{\usepackage{xurl}}{} % add URL line breaks if available
\IfFileExists{bookmark.sty}{\usepackage{bookmark}}{\usepackage{hyperref}}
\hypersetup{
  pdftitle={Enterprise Logging with Palo Alto Networks},
  pdfauthor={Franklin Diaz fdiaz@paloaltonetworks.com},
  pdflang={en},
  pdfkeywords={logging, syslog, tls},
  hidelinks,
  pdfcreator={LaTeX via pandoc}}
\urlstyle{same} % disable monospaced font for URLs
\usepackage{color}
\usepackage{fancyvrb}
\newcommand{\VerbBar}{|}
\newcommand{\VERB}{\Verb[commandchars=\\\{\}]}
\DefineVerbatimEnvironment{Highlighting}{Verbatim}{commandchars=\\\{\}}
% Add ',fontsize=\small' for more characters per line
\newenvironment{Shaded}{}{}
\newcommand{\AlertTok}[1]{\textcolor[rgb]{1.00,0.00,0.00}{\textbf{#1}}}
\newcommand{\AnnotationTok}[1]{\textcolor[rgb]{0.38,0.63,0.69}{\textbf{\textit{#1}}}}
\newcommand{\AttributeTok}[1]{\textcolor[rgb]{0.49,0.56,0.16}{#1}}
\newcommand{\BaseNTok}[1]{\textcolor[rgb]{0.25,0.63,0.44}{#1}}
\newcommand{\BuiltInTok}[1]{#1}
\newcommand{\CharTok}[1]{\textcolor[rgb]{0.25,0.44,0.63}{#1}}
\newcommand{\CommentTok}[1]{\textcolor[rgb]{0.38,0.63,0.69}{\textit{#1}}}
\newcommand{\CommentVarTok}[1]{\textcolor[rgb]{0.38,0.63,0.69}{\textbf{\textit{#1}}}}
\newcommand{\ConstantTok}[1]{\textcolor[rgb]{0.53,0.00,0.00}{#1}}
\newcommand{\ControlFlowTok}[1]{\textcolor[rgb]{0.00,0.44,0.13}{\textbf{#1}}}
\newcommand{\DataTypeTok}[1]{\textcolor[rgb]{0.56,0.13,0.00}{#1}}
\newcommand{\DecValTok}[1]{\textcolor[rgb]{0.25,0.63,0.44}{#1}}
\newcommand{\DocumentationTok}[1]{\textcolor[rgb]{0.73,0.13,0.13}{\textit{#1}}}
\newcommand{\ErrorTok}[1]{\textcolor[rgb]{1.00,0.00,0.00}{\textbf{#1}}}
\newcommand{\ExtensionTok}[1]{#1}
\newcommand{\FloatTok}[1]{\textcolor[rgb]{0.25,0.63,0.44}{#1}}
\newcommand{\FunctionTok}[1]{\textcolor[rgb]{0.02,0.16,0.49}{#1}}
\newcommand{\ImportTok}[1]{#1}
\newcommand{\InformationTok}[1]{\textcolor[rgb]{0.38,0.63,0.69}{\textbf{\textit{#1}}}}
\newcommand{\KeywordTok}[1]{\textcolor[rgb]{0.00,0.44,0.13}{\textbf{#1}}}
\newcommand{\NormalTok}[1]{#1}
\newcommand{\OperatorTok}[1]{\textcolor[rgb]{0.40,0.40,0.40}{#1}}
\newcommand{\OtherTok}[1]{\textcolor[rgb]{0.00,0.44,0.13}{#1}}
\newcommand{\PreprocessorTok}[1]{\textcolor[rgb]{0.74,0.48,0.00}{#1}}
\newcommand{\RegionMarkerTok}[1]{#1}
\newcommand{\SpecialCharTok}[1]{\textcolor[rgb]{0.25,0.44,0.63}{#1}}
\newcommand{\SpecialStringTok}[1]{\textcolor[rgb]{0.73,0.40,0.53}{#1}}
\newcommand{\StringTok}[1]{\textcolor[rgb]{0.25,0.44,0.63}{#1}}
\newcommand{\VariableTok}[1]{\textcolor[rgb]{0.10,0.09,0.49}{#1}}
\newcommand{\VerbatimStringTok}[1]{\textcolor[rgb]{0.25,0.44,0.63}{#1}}
\newcommand{\WarningTok}[1]{\textcolor[rgb]{0.38,0.63,0.69}{\textbf{\textit{#1}}}}
\usepackage{graphicx}
\makeatletter
\def\maxwidth{\ifdim\Gin@nat@width>\linewidth\linewidth\else\Gin@nat@width\fi}
\def\maxheight{\ifdim\Gin@nat@height>\textheight\textheight\else\Gin@nat@height\fi}
\makeatother
% Scale images if necessary, so that they will not overflow the page
% margins by default, and it is still possible to overwrite the defaults
% using explicit options in \includegraphics[width, height, ...]{}
\setkeys{Gin}{width=\maxwidth,height=\maxheight,keepaspectratio}
% Set default figure placement to htbp
\makeatletter
\def\fps@figure{htbp}
\makeatother
\setlength{\emergencystretch}{3em} % prevent overfull lines
\providecommand{\tightlist}{%
  \setlength{\itemsep}{0pt}\setlength{\parskip}{0pt}}
\setcounter{secnumdepth}{-\maxdimen} % remove section numbering
\ifxetex
  % Load polyglossia as late as possible: uses bidi with RTL langages (e.g. Hebrew, Arabic)
  \usepackage{polyglossia}
  \setmainlanguage[]{english}
\else
  \usepackage[shorthands=off,main=english]{babel}
\fi

\title{Enterprise Logging with Palo Alto Networks}
\usepackage{etoolbox}
\makeatletter
\providecommand{\subtitle}[1]{% add subtitle to \maketitle
  \apptocmd{\@title}{\par {\large #1 \par}}{}{}
}
\makeatother
\subtitle{Customer Logging Configuration}
\author{Franklin Diaz
\href{mailto:fdiaz@paloaltonetworks.com}{\nolinkurl{fdiaz@paloaltonetworks.com}}}
\date{March 24, 2025}

\begin{document}
\maketitle
\begin{abstract}
Instructions on configuration of TLS for an important account.
\end{abstract}

{
\setcounter{tocdepth}{3}
\tableofcontents
}
\hypertarget{introduction}{%
\section{Introduction}\label{introduction}}

\begin{figure}
\centering
\includegraphics{../images/customer/diagram1-logging.png}
\caption{an image showing the logging infra}
\end{figure}

\hypertarget{prerequisites}{%
\subsection{Prerequisites}\label{prerequisites}}

\begin{itemize}
\tightlist
\item
  Make available the TCP ports as shown in the document
  \href{https://docs.paloaltonetworks.com/strata-logging-service/administration/planning/ports-and-fqdns}{Strata
  Logging Service}
\item
  Use the support portal to create a Device Certificate on all Panorama
  and Firewall devices, and have Telemetry enabled.
\end{itemize}

\hypertarget{import-tls-certificates}{%
\subsection{Import TLS Certificates}\label{import-tls-certificates}}

The TLS certificates must be provided by the customer. Pleas take great
care when handling sensitive files such as keys, pem files,
certificates, etc. It is good practice to set ownership to
owner/read-only at all times. For exmaple \texttt{chmod\ 400\ *.pem}.

\begin{itemize}
\tightlist
\item
  For each desired service,
  \href{https://docs.paloaltonetworks.com/pan-os/11-1/pan-os-admin/certificate-management/obtain-certificates/import-a-certificate-and-private-key\#ide2b67a64-7100-4fa3-8304-48dbe64bcaa4}{generate
  or import a certificate on the firewall}
\item
  You can also use OCSP like so:
  \href{https://docs.paloaltonetworks.com/pan-os/11-1/pan-os-admin/certificate-management/obtain-certificates/obtain-a-certificate-from-an-external-ca}{Obtain
  a Certificate from an External CA}
\end{itemize}

\hypertarget{configure-ssltls-service-profile}{%
\subsection{Configure SSL/TLS Service
Profile}\label{configure-ssltls-service-profile}}

Making a service profile on the Panorama is the first step.

\begin{itemize}
\tightlist
\item
  \href{https://docs.paloaltonetworks.com/pan-os/11-1/pan-os-admin/certificate-management/configure-an-ssltls-service-profile/configure-an-ssltls-service-profile-pan-os}{Configure
  an SSL/TLS Service Profile (PAN-OS \& Panorama)}
\item
  Navigate to \texttt{Panorama} -\textgreater{}
  \texttt{Certificate\ Management} -\textgreater{}
  \texttt{SSL/TLS\ Service\ Profile}

  \begin{itemize}
  \tightlist
  \item
    set min TLS version to 1.2, max to 1.3
  \item
    Reference the Certificate you imported/created for logging.
  \end{itemize}
\end{itemize}

\hypertarget{configure-syslog-server-profile}{%
\subsection{Configure Syslog Server
Profile}\label{configure-syslog-server-profile}}

This step is only needed for SYSLOG configuration. The details about
each remote SYSLOG server are captured in these steps.

\begin{itemize}
\tightlist
\item
  Navigate to \texttt{Panorama} -\textgreater{}
  \texttt{Server\ Profiles} -\textgreater{} \texttt{SYSLOG} and create
  each logging profile as desired.

  \begin{itemize}
  \tightlist
  \item
    The TCP profile" for splunk testing is server 10.228.96.182,
    transport TCP, port 8642, format IETF
  \item
    The TLS profile" for splunk testing is server 10.228.96.182,
    transport TLS, port 8643. format IETF
  \end{itemize}
\item
  We do not usually change any settings on the
  \texttt{Custom\ Log\ Format} tab.
\end{itemize}

\hypertarget{configure-log-settings}{%
\subsection{Configure Log Settings}\label{configure-log-settings}}

\begin{itemize}
\tightlist
\item
  Navigate to \texttt{Panorama} -\textgreater{} \texttt{Log\ Settings}
  -\textgreater{} \texttt{System}
\item
  Add the new SYSLOG profiles here with the \texttt{All\ Logs} filter.
\end{itemize}

\hypertarget{qradar}{%
\subsubsection{QRadar}\label{qradar}}

The QRadar server is at \texttt{qrocvip-mah.xxx.com} with IP
\texttt{153.2.226.136} on TCP port 514.

\hypertarget{syslog-client-configuration}{%
\subsubsection{Syslog Client
Configuration}\label{syslog-client-configuration}}

Syslog is configured in two places on the Panorama. Once in the Panorama
tab for the Panorama itself, and again in the Device Templates for the
Firewall devices.

To configure firewall SYSLOG from the Panorama console:

\begin{itemize}
\tightlist
\item
  From the Global Template:
  \texttt{device\ -\textgreater{}\ server\ profiles\ -\textgreater{}\ syslog}

  \begin{itemize}
  \tightlist
  \item
    This makes the firewalls aware of all the SYSLOG servers.
  \end{itemize}
\item
  For SYSLOG coming from the Panorama(s) you should also check
  \texttt{Panorama\ -\textgreater{}\ server\ profiles\ -\textgreater{}\ syslog}
\item
  NOTE: SYSLOG is system configuration logging, which is different from
  the traffic log forwarding.
\item
  Make sure the IP address of the jump host or other test host is in the
  permitted mgmt IP address list.
\item
  You could also try tcpdump from
  \texttt{Panorama\ \textgreater{}\ Managed\ Devices\ \textgreater{}\ Troubleshooting}
\end{itemize}

Validate the basic connectivity between Panorama and SYSLOG server.
These example commands are meant to be run from the CLI.

\begin{Shaded}
\begin{Highlighting}[]
\ExtensionTok{debug}\NormalTok{ syslog{-}params show}
\FunctionTok{ping}\NormalTok{ host }\OperatorTok{\textless{}}\NormalTok{IP address of syslog server}\OperatorTok{\textgreater{}}
\ExtensionTok{traceroute}\NormalTok{ host }\OperatorTok{\textless{}}\NormalTok{IP address of syslog server}\OperatorTok{\textgreater{}}
\end{Highlighting}
\end{Shaded}

\begin{itemize}
\tightlist
\item
  Check the TCP connection between the firewall and the syslog server.
\item
  Replace TCP port if not using TCP/514, for exmaple 8642 or 8643.
\end{itemize}

\begin{Shaded}
\begin{Highlighting}[]
\ExtensionTok{show}\NormalTok{ netstat numeric{-}hosts yes numeric{-}ports yes all yes }\KeywordTok{|} \ExtensionTok{match}\NormalTok{ 514}
\ExtensionTok{show}\NormalTok{ netstat numeric{-}hosts yes numeric{-}ports yes all yes }\KeywordTok{|} \ExtensionTok{match}\NormalTok{ 8642}
\end{Highlighting}
\end{Shaded}

\begin{itemize}
\tightlist
\item
  On Panorama Template, check permitted IP addresses

  \begin{itemize}
  \tightlist
  \item
    \texttt{Device\ \textgreater{}\ Setup\ \textgreater{}\ Interfaces\ \textgreater{}\ Management\ \textgreater{}\ Permitted\ IP\ Addresses}
  \item
    The IP address for each SYSLOG collector should be listed in the
    permitted list.
  \end{itemize}
\end{itemize}

\hypertarget{cortex-data-lake}{%
\subsubsection{Cortex Data Lake}\label{cortex-data-lake}}

\begin{itemize}
\tightlist
\item
  Capture and describe the settings for CDL here.
\item
  Note that CDL does not work without a valid device certificate on the
  host.
\end{itemize}

\end{document}
